\documentclass[11pt]{article}
\usepackage[margin=1in]{geometry}
\usepackage{booktabs}
\usepackage{graphicx}
\usepackage{float}
\usepackage{caption}
\usepackage{hyperref}
\usepackage[numbers,sort&compress]{natbib}  % Sage Vancouver: numbered
\usepackage{amsmath}
\usepackage{array}

\captionsetup{font=small,labelfont=bf}
\graphicspath{{../05_Figures/}{../08_Model/}}
\makeatletter
\def\input@path{{../06_Tables/}{../08_Model/}}
\makeatother

\title{A Bayesian Left-Censored Joint Longitudinal--Interval Model for
Irregular Molecular Monitoring in Chronic Myeloid Leukemia}

%% ---- TITLE PAGE INFORMATION (submit as a separate, non-anonymized file) ----
%% Xulin Pan (corresponding author), xxp5066@psu.edu
%%   Pennsylvania State University, University Park, PA 16802, USA
%% Jun Lv, School of Mathematics and Statistics, Yunnan University, Kunming, China
%% Chuanming Wang, School of Biological and Food Engineering, Honghe University, Mengzi, China
%% Xueren Wang, School of Mathematics and Statistics, Yunnan University, Kunming, China
%% ---------------------------------------------------------------------------
\author{}     % anonymized for double-anonymized peer review
\date{}

\graphicspath{{../05_Figures/}{../08_Model/}}
\def\input@path{{../06_Tables/}{../08_Model/}}
\renewcommand{\thesection}{S\arabic{section}}
\renewcommand{\thetable}{S\arabic{table}}
\renewcommand{\thefigure}{S\arabic{figure}}

\title{Supplementary Material\\[2pt]
\large A Bayesian Joint Longitudinal--Interval and Multi-State
Threshold-Crossing Framework for Irregular Molecular Monitoring in Chronic
Myeloid Leukemia}
\author{}
\date{}

\begin{document}
\maketitle

\noindent This supplement supports the main paper with: (i) the closed-form
running extrema and the optional relapse frailty of the multi-state
threshold-crossing model (Section~3.3 of the main paper); (ii) an exploratory
informative visit-process model; and (iii) two secondary simulation scenarios
(endpoint choice and informative monitoring). The multi-state model, its
simulation recovery, and its application to the cohort are reported in the main
paper.

\section{Closed-form running extrema and relapse frailty}
\label{sec:s-frailty}

Because the latent trajectory $m_i(\ell)$ is quadratic in $\ell=\log(1+t)$
(piecewise-quadratic once a spline slope change is added), the running extrema
that define the multi-state transitions are available in closed form. The
running minimum equals the value at the vertex
$\ell^\ast=-(\beta_1+b_{1i})/(2\beta_2)$ clamped to the observed range
$[0,\ell(t)]$; past the vertex $m_i$ is monotone increasing, so the upward
(relapse) crossing and the durability maximum are point evaluations. This
removes the soft-minimum grid approximation entirely and yields a smooth log
density, which in practice eliminates the divergent transitions the grid
induces.

Between-patient heterogeneity in relapse propensity can be captured by a
Gaussian frailty $\rho_i\sim N(0,\sigma_\rho^2)$ on the relapse threshold.
Because
\[
\int\Phi\!\left\{\frac{x+\rho}{\sigma_{\mathrm{thr}}}\right\}\,\mathrm dN(\rho;0,\sigma_\rho^2)
=\Phi\!\left\{\frac{x}{\sqrt{\sigma_{\mathrm{thr}}^2+\sigma_\rho^2}}\right\},
\]
the frailty is equivalent to a wider relapse-specific width
$\sigma_{\mathrm{rel}}=\sqrt{\sigma_{\mathrm{thr}}^2+\sigma_\rho^2}$ and
preserves the closed form. In the cohort the frailty modestly improved relapse
calibration but was only weakly identified (9 confirmed relapses) and added
little over the spline model in simulations where trajectory curvature already
explains relapse; the base spline model is therefore preferred.

\section{Exploratory informative visit-process model (not fitted)}
\label{sec:s-visit}

Monitoring intensity may itself depend on molecular response: a patient doing
well may be seen less often, and more frequent testing yields more
opportunities to document DMR. This dependence is written into the governing
guideline, which raises molecular-monitoring frequency when the BCR-ABL
transcript rises \cite{CSH2011CMLGuideline}. To assess sensitivity to
informative visiting, one may model the visit intensity jointly with the
trajectory,
\[
\lambda_i^V(t)=
\exp\{\delta_0+\delta_1[m_i(t)+4.5]+\delta_2\log(1+t)+u_i\},
\qquad u_i\sim N(0,\sigma_u^2),
\]
so that visit intensity depends on the latent MRD, time, and a patient-level
monitoring random effect $u_i$. For visit times $v_{i1},\ldots,v_{iM_i}$ in
$[E_i,C_i]$ the visit-process likelihood is
\[
L_i^V=
\left\{\prod_{\ell=1}^{M_i}\lambda_i^V(v_{i\ell})\right\}
\exp\left\{-\int_{E_i}^{C_i}\lambda_i^V(s)\,ds\right\},
\]
with the cumulative intensity evaluated by quadrature. Because this extension
adds weakly identified parameters in a cohort of 87 patients, it is presented
as a sensitivity model to be evaluated primarily by simulation and in larger
cohorts.

\section{Simulation: endpoint choice and informative monitoring}
\label{sec:s-sim}

The main-text simulation focuses on the primary joint model and the multi-state
recovery study. Here we report the two secondary scenarios summarised in the
patient-level calibration table (Table~12 of the main paper).

\emph{Endpoint choice.} A latent threshold-crossing estimator, which targets
biological onset rather than the visit-documented endpoint, under-predicted
documented DMR more strongly than the floor model (mean predicted $\approx0.44$
versus observed $\approx0.85$ at $n=150$; difference $\approx-0.41$, versus
$\approx-0.26$ for the floor model). This quantifies the distinction between the
latent-crossing and documented-DMR estimands: because documentation occurs only
at visits, an estimator that targets biological onset necessarily predicts fewer
\emph{documented} events under interval ascertainment.

\emph{Informative monitoring.} Under a response-adaptive visit process the
observed documented-DMR rate fell (to $\approx0.59$) and the floor model's
calibration slope rose above one ($1.16$), showing that the monitoring process
materially affects calibration. This motivates the joint visit-process
extension outlined as future work in the main paper.

\bibliographystyle{unsrtnat}
\bibliography{../07_References/peer_reviewed_papers_selected}

\end{document}
